\chapter{Conclusion} \label{Chap8}

\section{Summary and Contribution}

This dissertation set out, against the five objectives stated in Chapter~\ref{Chap1}, to develop
and evaluate machine learning approaches for CH\textsubscript{4} forecasting at the North Wyke Farm
Platform and to demonstrate their integration with scenario projection under contrasting management
and climate conditions. Each objective has been met directly: a systematic review of the relevant
literature (Chapter~\ref{Chap2}) established that ecosystem-scale CH\textsubscript{4} forecasting at
a managed temperate grassland was, to the author's knowledge, previously unaddressed; NWFP's EC
CH\textsubscript{4} record was acquired, corrected, and gap-filled to a continuous series
(Chapter~\ref{Chap3}, Chapter~\ref{Chap4}); a benchmarking pipeline spanning statistical baselines,
tree ensembles, and tabular foundation models was designed and evaluated under temporal
cross-validation throughout (Chapter~\ref{Chap4}, Chapter~\ref{Chap5}); interpretability and
uncertainty quantification were applied across every stage (Chapter~\ref{Chap4},
Chapter~\ref{Chap5}, Chapter~\ref{Chap6}); and synthetic management and climate scenarios were
generated and evaluated against the resulting models (Chapter~\ref{Chap6}).

The research questions posed in Chapter~\ref{Chap1} are each answered directly by this work. RQ1's
survey of ML approaches for CH4 gap-filling and its distinction from multi-step forecasting is
answered in Chapter~\ref{Chap2}: gap-filling and forecasting are structurally different problems,
and forecasting at this ecosystem scale had no prior treatment in the literature. RQ2's comparison
of statistical baselines, tree ensembles, and deep learning under NWFP's own temporal variability is
answered in Chapters~\ref{Chap4} and~\ref{Chap5}: tabular foundation models lead, tree ensembles
follow closely, and deep sequence models trail, for the two separable reasons established in
Chapter~\ref{Chap7}. RQ3's question of transferability from adjacent domains is answered throughout
Chapter~\ref{Chap4}: some published findings (Random Forest beating MDS, ANN's competitiveness at
longer gaps) replicate cleanly at NWFP, while others (PCA degrading ML performance) do not,
underscoring that findings from wetland and rice-paddy sites do not transfer uniformly to managed
UK grassland. RQ4's question of how interpretability and uncertainty quantification decode the
interaction between continuous environmental drivers and discrete management events is answered by
this dissertation's single most robust finding, discussed below. RQ5's question of the structural
requirements for integrating forecasting into scenario-based "what-if" analysis is answered by
Chapter~\ref{Chap6}'s working scenario architecture and, equally importantly, by the structural
limitations that architecture surfaces (Section 7.2) -- what such a system requires is now
considerably better understood than it was at the outset of this project, including where it
currently falls short.

If this dissertation is remembered for a single finding, it should be this one: \textbf{livestock
density is the dominant driver of CH\textsubscript{4} flux at NWFP, and this is not a fragile,
single-method result.} It recurs independently across gap-filling feature importance, forecasting
SHAP attribution and permutation importance, SARIMAX regression coefficients, and scenario
projection's own livestock-ladder response -- genuinely different methods, different model
architectures, and different chapters of this dissertation, each arriving at the same answer
independently. Most empirical work produces one headline number and hopes it survives scrutiny; this
project instead produced a finding that shows up no matter which door is used to go looking for it.

Alongside the modelling results, this dissertation's evaluation-methodology corrections deserve
explicit credit rather than being read as small next to the headline model numbers, since several of
them were larger in effect than any single feature-engineering gain in the entire project. Tower 2's
evaluation-protocol fix alone (Section 4.4.3) moved its gap-filling sklearn $R^2$ from $-16.9$ under
an early, seasonally-mismatched train/test split to a positive result under the corrected
full-period gap-CV protocol used throughout Chapter~\ref{Chap4}'s improved pipeline
(Section 4.5.3.1) -- a bigger single
correction than any feature added across the whole gap-filling programme. The reversal from
single-anchor to multi-anchor rollout evaluation is a genuine, non-obvious methodological lesson in
its own right, not a footnote: single-anchor rankings in this project's own preliminary work proved
unreliable and would have supported different conclusions had they been trusted. And catching the
observed-versus-gap-filled-target scoring circularity (Section 5.1.3) before it was allowed to
corrupt a headline forecasting claim is exactly the kind of methodological rigour that is easy to
skip and expensive to get wrong.

Finally, it is worth being precise about what this dissertation has, and has not, shown. On
gap-filling, this project has demonstrated a significant, benchmarked improvement over traditional
methods such as MDS and plain RFm, achieved through a combination of tabular foundation models and a
substantially richer feature set than the published baselines this work replicates. On forecasting,
this project has generated models that reliably beat naive and climatological baselines, but does
not yet produce fully trustworthy point forecasts, particularly during the high-magnitude spike
events that dominate CH\textsubscript{4}'s error budget throughout this dissertation. Scenario
projection, in turn, demonstrates that the resulting architecture can produce distinguishable,
interpretable projections under contrasting management and climate conditions, while carrying the
structural limitation that no future scenario can ever be checked against ground truth. Taken
together, this dissertation delivers a benchmarked, uncertainty-quantified, and interpretable
foundation for CH\textsubscript{4} forecasting and scenario projection at a managed UK grassland
site that did not previously exist in the literature -- not a finished, deployment-ready forecasting
system, but a substantial and honestly-bounded step toward one.
